\section{The equality case gives a simplicial cycle}
\label{sec:topology}

The symmetric-matrix extension gives more than a lower bound on the
number of labels.  At equality, its degree controls every maximal
support of the admissible map.  We prove that these supports all have
the same size and that their sum is a cycle.  No triangulation of
projective space is assumed.

Throughout this section, homology and cohomology have coefficients in
$\mathbb F_2$.  For a Euclidean space $V$ of
dimension $k$, write
\[
 \Sigma(V)=\{A\in\operatorname{Sym}(V):\operatorname{tr}|A|=1\},
 \qquad
 \Sigma_m=\{y\in\mathbb R^m:\textstyle\sum_i|y_i|=1\}.
\]
Radial projection identifies these spaces antipodally with
$S^{N_k-1}$ and $S^{m-1}$, respectively.  We give them the smooth
structures transported by these identifications.  Near nonsingular
operators, respectively full-support label vectors, these structures
agree with their usual smooth hypersurface structures: the relevant
norm functions are smooth there.  Let
$f:\mathbb P(V)\to\Delta^{m-1}$ be a smooth admissible map, and let
$F:\Sigma(V)\to\Sigma_m$ be its continuous odd extension constructed
in the preceding section.

\begin{proposition}[The number of labels]\label{thm:label-bound}
If $k\ge2$, every smooth admissible map
$f:\mathbb P(\mathbb R^k)\to\Delta^{m-1}$ satisfies $m\ge N_k$.
If $m=N_k$, its extension $F$ has mod-two degree one and is
surjective.  In this equality case, every pair of labels is an edge
of the support complex of $f$.
\end{proposition}

\begin{proof}
An odd map $S^n\to S^d$ requires $n\le d$, and an odd self-map
of $S^n$ has odd degree; see \cite[Proposition~2B.6 and
Corollary~2B.7]{Hatcher2002}.  Applying these facts to $F$ proves the
bound and the degree assertion.  A map of nonzero degree cannot omit
a point of its target sphere, since the complement of a point is
contractible.

Suppose $m=N_k$ and fix distinct labels $i,j$.  Choose a vector in
$\Sigma_m$ with every coordinate nonzero, with precisely the
$i$th and $j$th coordinates positive.  At a preimage $A$ under $F$,
every line in the positive spectral subspace of $A$ uses only these
two labels.  That subspace cannot have dimension at least two:
restriction to a two-dimensional subspace would contradict the bound
$N_2=3$ just proved.  Its dimension is therefore one.  The exact
spectral-support property of $F$ says that this line has both
labels $i,j$ in its support, so $\{i,j\}$ is an edge.
\end{proof}

We will count regular preimages of maps which need not be smooth
everywhere.  The following local form of the degree formula explains
why this causes no problem.

\begin{lemma}[Local counting]\label{lem:local-degree}
Let $M$ be a closed smooth $n$-manifold, possibly disconnected, and
let $q:M\to Y$ be continuous, with $Y$ Hausdorff.  Suppose a point $y\in Y$ has a
neighborhood homeomorphic to an open subset of $\mathbb R^n$, and
that $q$ is smooth in the corresponding coordinates near
$q^{-1}(y)$, with invertible derivative at each point of that fiber.
Then the fiber is finite, and the image of the total fundamental
class of $M$ in
\[
 H_n(Y,Y\setminus\{y\})\cong\mathbb F_2
\]
is $|q^{-1}(y)|$ modulo two.  In particular, if $Y$ is a connected
closed $n$-manifold, this parity is the mod-two degree of $q$.
\end{lemma}

\begin{proof}
The inverse function theorem makes every point of the fiber
isolated.  The fiber is closed in compact $M$, hence finite.
By excision,
\[
 H_n(M,M\setminus q^{-1}(y))
 \cong\bigoplus_{x\in q^{-1}(y)}H_n(M,M\setminus\{x\}).
\]
The fundamental class has coordinate one in each summand.  A local
diffeomorphism sends each local generator to the generator at $y$.
Naturality of relative homology now gives the assertion.  This
argument uses no smoothness away from the selected fiber.
\end{proof}

\begin{theorem}[Supports at equality]\label{thm:equality-complex}
Let $k\ge3$ and $m=N_k$.  Suppose
$f:\mathbb P(\mathbb R^k)\to\Delta^{m-1}$ is a smooth admissible map.
Its support complex $K$ has every pair of labels as an edge, every
facet has $k$ vertices, and
\[
 \sum_{I\text{ a facet of }K} I
\]
is a simplicial cycle over $\mathbb F_2$.  Consequently every
$(k-2)$-dimensional face lies in a positive even number of facets.
\end{theorem}

\begin{proof}
The edge assertion is Proposition~\ref{thm:label-bound}.  We prove
the remaining assertions by examining an arbitrary facet $I$.
Set
\[
 s=|I|,\qquad J=\{1,\ldots,m\}\setminus I,\qquad
 a=|J|,\qquad h=N_{k-1}.
\]
Since $I$ is a maximal support face, there is a line $x_0$ with
$S(x_0)=I$.  Every line in $x_0^\perp$ uses only labels in $J$.
Proposition~\ref{thm:label-bound}, applied in $x_0^\perp$, gives
\begin{equation}\label{eq:facet-size-upper}
 a\ge h,\qquad s=m-a\le N_k-N_{k-1}=k.
\end{equation}

\paragraph{A compact level set over one affine chart.}
Define
\[
 U_I=\{x\in\mathbb P(V):f_i(x)>0\text{ for every }i\in I\}.
\]
Maximality of $I$ implies $S(x)=I$ for all $x\in U_I$.  Thus the
coordinate map $f_I$ takes $U_I$ smoothly to the open simplex
$\mathring\Delta_I$ on $I$.  Moreover, $U_I$ lies in the affine
projective chart
\[
 \mathbb P(V)\setminus\mathbb P(x_0^\perp),
\]
because a line orthogonal to $x_0$ cannot use any label in $I$.
Choose a regular value $r\in\mathring\Delta_I$ of $f_I$, using
Sard's Theorem \cite[Theorem~6.10]{Lee2013}, and put
\[
 L=f_I^{-1}(r)\subset U_I.
\]
For now $L$ is allowed to be empty.  If $\widetilde r$ denotes $r$
extended by zero on $J$, then
$L=f^{-1}(\widetilde r)$: positivity of the coordinates of $r$
puts every such preimage in $U_I$.  Hence $L$ is compact.  It is a
smooth manifold without boundary \cite[Corollary~5.14]{Lee2013} of dimension
\begin{equation}\label{eq:level-dimension}
 \dim L=(k-1)-(s-1)=k-s=a-h.
\end{equation}

If $s<k$, the level-set dimension in \eqref{eq:level-dimension} is positive.
We will construct a map on a sphere bundle over $L$ whose regular
preimages have odd cardinality, by comparison with $F$. The affine-chart
trivialization will force the induced map on the projective bundle to
have degree zero when $\dim L>0$, contradicting that count. When $s=k$,
the same count will determine the coefficient of $I$ in the desired cycle.

For $x\in U_I$, admissibility ensures that $f$ on
$\mathbb P(x^\perp)$ uses only labels in $J$.  Apply the extension
construction in this $(k-1)$-dimensional space to obtain an odd map
\[
 G_x:\Sigma(x^\perp)\longrightarrow\Sigma_a.
\]
The subspace-restriction and continuity properties of the extension
make these maps a continuous map on the corresponding sphere bundle.
Restrict that bundle to $L$ and write
\[
 M=\{(x,B):x\in L,\ B\in\operatorname{Sym}(x^\perp),
                     \ \operatorname{tr}|B|=1\},
 \qquad G(x,B)=G_x(B).
\]

The bundle $x\mapsto x^\perp$ is trivial on the affine chart above.
Indeed, choose the unit representative $u$ of $x$ having positive
inner product with a fixed unit representative of $x_0$.  Orthogonal
projection from $x_0^\perp$ onto $u^\perp$ is invertible and varies
smoothly with $x$; orthonormalizing its images of a fixed basis gives
a smooth orthonormal trivialization.  The induced trivialization
of symmetric operators, followed by radial change from trace norm
to a Euclidean norm, gives
\begin{equation}\label{eq:product-sphere-bundle}
 M\cong L\times S^{h-1},\qquad
 (x,B)\longmapsto(x,-B)\quad\leftrightarrow\quad(x,z)\longmapsto(x,-z).
\end{equation}
Equip $M$ with the smooth structure given by the Euclidean sphere
bundle.  The trace-norm parametrization is smooth near every
definite $B$, which is all that will be needed.  In particular,
\begin{equation}\label{eq:bundle-dimension}
 \dim M=(a-h)+(h-1)=a-1.
\end{equation}

\paragraph{Odd parity on the bundle.}
Consider the open negative orthant of $\Sigma_a$, where all $J$
coordinates are strictly negative.  At a preimage of this orthant,
$B$ has no positive spectral subspace.  It has no kernel either:
the exact support property says that all $J$ labels occur on its
range, leaving no label for a kernel line.  Thus $B$ is negative
definite on $x^\perp$.  Consequently $G$ is smooth on the preimage
of the negative orthant.  Sard's Theorem allows us to choose a
regular value $g$ in that orthant.  Its preimage is finite by
compactness, including the possibility of an empty preimage.

Fix $0<\lambda<1$ and form the full-support vector
\[
 y=(\lambda r,(1-\lambda)g)\in\Sigma_m,
\]
with the two coordinate blocks indexed by $I$ and $J$.  We claim
that $F^{-1}(y)$ corresponds exactly to $G^{-1}(g)$.
Every $A\in F^{-1}(y)$ is nonsingular by the full-support property.
All lines in its nonzero positive spectral subspace use only
labels in $I$.  That subspace must be one-dimensional: a subspace
of dimension at least two intersects the hyperplane $x_0^\perp$,
where no label in $I$ is available.  Call the resulting positive
line $x$.  The positive output block has every $I$ coordinate
nonzero, so $S(x)=I$ and $x\in U_I$.  Since the sum of the positive
coordinates of $F(A)$ is $\lambda$, the positive eigenvalue is
$\lambda$.  Therefore
\begin{equation}\label{eq:preimage-splitting}
 A=\lambda P_x+(1-\lambda)B,
 \qquad B<0\text{ on }x^\perp,\quad \operatorname{tr}|B|=1.
\end{equation}
The subspace formula for the extension now gives
$f_I(x)=r$ and $G_x(B)=g$.  Conversely, every $(x,B)\in G^{-1}(g)$
gives a preimage by \eqref{eq:preimage-splitting}.

The value $y$ is regular in the local sense of
Lemma~\ref{lem:local-degree}.  To see this precisely, near a
preimage use $(\lambda',x',B')$ as coordinates for symmetric
operators of signature $(1,k-1)$, with $B'$ negative definite and
normalized on $x'^\perp$.  The single positive eigenvalue and its
line vary smoothly, so these are smooth coordinates on the
trace-norm sphere near that preimage.  In the target orthant, use
the positive mass followed by the two normalized coordinate
blocks.  In these coordinates $F$ is
\begin{equation}\label{eq:triangular-derivative}
 (\lambda',x',B')\longmapsto
       (\lambda',f_I(x'),G_{x'}(B')).
\end{equation}
The first derivative block is the identity.  The second is
surjective because $r$ is regular for $f_I$.  On the kernel of
these two blocks, the third block is the derivative of
$G|_M$, which is surjective because $g$ is regular.  Thus the full
derivative is surjective, and is an isomorphism since source and
target both have dimension $m-1$.  By the odd degree of $F$ and
Lemma~\ref{lem:local-degree},
\begin{equation}\label{eq:bundle-odd-count}
 |G^{-1}(g)|=|F^{-1}(y)|\equiv1\pmod2.
\end{equation}
This proves, in particular, that $L$ is nonempty.

\paragraph{The double cover forces maximal supports to have size $k$.}
The oddness of $G$ in the sphere variable gives a quotient map
\[
 \overline G:L\times\mathbb{RP}^{h-1}
                  \longrightarrow\mathbb{RP}^{a-1}.
\]
Write $w\in H^1(\mathbb{RP}^{a-1})$ and
$u\in H^1(\mathbb{RP}^{h-1})$ for the first Stiefel--Whitney classes
of the real line bundles associated with their antipodal double covers.
Equivariance identifies the pullback of the target
sphere cover with
\[
 L\times S^{h-1}\longrightarrow L\times\mathbb{RP}^{h-1}.
\]
Indeed, the identification sends $(x,z)$ to
$((x,[z]),G(x,z))$ and is an isomorphism on each two-point fiber.
This cover isomorphism induces an isomorphism of the associated real
line bundles. Their first Stiefel--Whitney classes are natural under
pullback \cite[Theorem~3.1(a)]{HatcherVB2017}.
The source cover is pulled back from the second factor, so
\begin{equation}\label{eq:cover-class}
 \overline G^*w=\operatorname{pr}_2^*u.
\end{equation}
There is no contribution from $H^1(L)$ in this identity.

Suppose $a>h$.  The projective-space cohomology calculation
\cite[Theorem~3.19]{Hatcher2002} gives
\[
 u^h=0,\qquad
 \overline G^*(w^{a-1})=
       \operatorname{pr}_2^*(u^{a-1})=0.
\]
Both source and target of $\overline G$ are closed manifolds of
dimension $a-1$, and $w^{a-1}$ is the target's top cohomology
generator.  Evaluation on the total fundamental class therefore
gives mod-two degree zero.  On the other hand, $[g]$ is a regular
value in local quotient coordinates, and every preimage of $[g]$
has exactly one representative mapping to $g$.  Its preimage count
is therefore the odd number in \eqref{eq:bundle-odd-count}, a
contradiction.  Together with \eqref{eq:facet-size-upper}, this
proves $a=h$ and $s=k$ for every facet.

\paragraph{The coefficients of the facet cycle.}
Now $\dim L=0$, so $L$ is a finite set.  For each $x\in L$, the
map $G_x$ is odd between spheres of the same dimension $h-1$.
Regularity of $g$ for $G$ implies its regularity for each such
fiber map.  Each fiber therefore contributes an odd number to
\eqref{eq:bundle-odd-count}, and hence
\begin{equation}\label{eq:facet-odd-count}
 |f^{-1}(\widetilde r)|=|L|\equiv1\pmod2.
\end{equation}

Every facet has $k$ vertices, so $K$ has dimension $k-1$.
Using the identification of simplicial and singular homology
\cite[Theorem~2.27]{Hatcher2002}, consider the homology class
\[
 f_*[\mathbb P(V)]\in H_{k-1}(|K|).
\]
The interior point $\widetilde r$ of the facet $I$ has a
neighborhood in $|K|$ contained in that facet.  Passing to
$H_{k-1}(|K|,|K|\setminus\{\widetilde r\})$ reads the coefficient
of $I$ in a simplicial top cycle.  By the regularity of $r$,
Lemma~\ref{lem:local-degree} and \eqref{eq:facet-odd-count}, this
coefficient is one.  The same holds for every facet.  Since a
$(k-1)$-dimensional complex has no simplicial $k$-chains, its
top homology classes are exactly its top cycles.  Thus
$f_*[\mathbb P(V)]$ is represented by the sum of all facets.
Its boundary vanishes.  Each codimension-one face consequently
occurs in an even number of facets, and purity makes that number
positive.
\end{proof}
